% Copyright 2026  Ed Bueler

\documentclass[10pt,
               svgnames,
               hyperref={colorlinks,
                         citecolor=DeepPink4,
                         linkcolor=FireBrick,
                         urlcolor=blue},
               usepdftitle=false]{beamer}

\usetheme{Madrid}
\usecolortheme{beaver}
\setbeamercovered{transparent}
\setbeamerfont{frametitle}{size=\large}
\setbeamercolor*{block title}{bg=red!10}
\setbeamercolor*{block body}{bg=red!5}

\usepackage[english]{babel}
\usepackage[latin1]{inputenc}
\usepackage{times}
\usepackage[T1]{fontenc}
% Or whatever. Note that the encoding and the font should match. If T1
% does not look nice, try deleting the line with the fontenc.

\usepackage{fancyvrb,empheq,bm,xspace,dsfont,esint}
\usepackage{minted}
\usepackage{tikz}
\usetikzlibrary{positioning,arrows.meta}

% If you wish to uncover everything in a step-wise fashion, uncomment
% the following command: 
%\beamerdefaultoverlayspecification{<+->}

\newcommand{\bb}{\mathbf{b}}
\newcommand{\bc}{\mathbf{c}}
\newcommand{\bbf}{\mathbf{f}}
\newcommand{\bl}{\bm{\ell}}
\newcommand{\bn}{\mathbf{n}}
\newcommand{\bq}{\mathbf{q}}
\newcommand{\br}{\mathbf{r}}
\newcommand{\bs}{\mathbf{s}}
\newcommand{\bx}{\mathbf{x}}
\newcommand{\by}{\mathbf{y}}
\newcommand{\bv}{\mathbf{v}}
\newcommand{\bu}{\mathbf{u}}
\newcommand{\bw}{\mathbf{w}}

\newcommand{\bE}{\mathbf{E}}
\newcommand{\bV}{\mathbf{V}}

\newcommand{\cB}{\mathcal{B}}
\newcommand{\cE}{\mathcal{E}}
\newcommand{\cM}{\mathcal{M}}
\newcommand{\cR}{\mathcal{R}}

\newcommand{\bzero}{\bm{0}}

\newcommand{\CC}{\mathbb{C}}
\newcommand{\NN}{\mathbb{N}}
\newcommand{\RR}{\mathbb{R}}
\newcommand{\ZZ}{\mathbb{Z}}

\newcommand{\one}{\mathds{1}}

\newcommand{\Div}{\nabla\cdot}
\newcommand{\grad}{\nabla}

\renewcommand{\Re}{\operatorname{Re}}
\renewcommand{\Im}{\operatorname{Im}}

\newcommand{\range}{\operatorname{range}}
\newcommand{\Span}{\operatorname{span}}

\newcommand{\Matlab}{\textsc{Matlab}\xspace}
\newcommand{\eps}{\epsilon}

\newcommand{\Lloc}{L_{\text{loc}}^1}

\newcommand{\ds}{\displaystyle}

\newcommand{\ip}[2]{\left<#1,#2\right>}

\newcommand{\twovect}[4]{\ensuremath{{#1}_{#2} =
                            \begin{bmatrix} #3 \\ #4 \end{bmatrix}}}

\newcommand{\rbullet}{{\color{FireBrick} \bullet}}

\newcommand*\circled[1]{\tikz[baseline=(char.base)]{
            \node[shape=circle,draw,inner sep=2pt] (char) {#1};}}

\newcommand{\textbook}{Saxe, \emph{Beginning Functional Analysis}}

\DefineVerbatimEnvironment{mVerb}{Verbatim}{numbersep=2mm,
frame=lines,framerule=0.1mm,framesep=2mm,xleftmargin=4mm,fontsize=\scriptsize}

\newcommand{\itemnum}[1]{\setcounter{enumi}{#1}\usebeamertemplate{enumerate item}}



\title[Poisson equation: theory and approximation]{Poisson equation: some theory, and approximation}

\subtitle{calculations for week 13}

\author{Ed Bueler}

\institute[]{UAF Math 617 Functional Analysis}

\date{Spring 2026}


\begin{document}
\beamertemplatenavigationsymbolsempty

\begin{frame}
  \maketitle
\end{frame}


\begin{frame}{Outline}
  \tableofcontents[hideallsubsections]
\end{frame}

\AtBeginSection[]
{% do not put toc here this time
}


\section{recall the Poisson equation}


\begin{frame}{recall: boundary-value problems for the Poisson equation}

\begin{itemize}
\item suppose $\Omega\subset \RR^d$ is open and has Lipschitz (or continuously-differentiable) boundary
\end{itemize}

\begin{block}{general Dirichlet problem (\emph{strong form})}
for $f\in L^2(\Omega)$ and $g_D\in C^1(\partial\Omega)$, find $u:\Omega\to \RR$ so that
\begin{align*}
- \grad^2 u &= f \quad \text{ on } \Omega \\
u &= g_D \quad \text{ on } \partial \Omega
\end{align*}
\end{block}

\begin{block}{general Neumann problem (\emph{strong form})}
for $f\in L^2(\Omega)$ and $g_N\in L^2(\partial\Omega)$, find $u:\Omega\to \RR$ so that
\begin{align*}
- \grad^2 u &= f \quad \text{ on } \Omega \\
\grad u\cdot \bn &= g_N \quad \text{ on } \partial \Omega
\end{align*}
where $\bn$ is the outward unit normal vector field along $\partial\Omega$
\end{block}

\begin{itemize}
\item other notation: $\grad u\cdot \bn = \frac{\partial u}{\partial n}$
\end{itemize}
\end{frame}


\begin{frame}{recall: the weak form}

\begin{block}{weak formulation of the Poisson equation}
Suppose $f\in L^2(\Omega)$ and $g_N\in L^2(\partial\Omega)$ are given.  Solve one of these problems:
\begin{enumerate}
\item Find a solution $u\in H_0^1(\Omega)$ so that
    $$\int_\Omega \grad u \cdot \grad v\,dx = \int_\Omega f v\,dx \hspace{20mm} \text{(homogeneous Dirichlet)}$$
for all test functions $v\in H_0^1(\Omega)$.
\item Find a solution $u\in H^1(\Omega)$ so that
    $$\int_\Omega \grad u \cdot \grad v\,dx = \int_\Omega f v\,dx + \int_{\partial\Omega} g_N v\,ds \hspace{10mm} \text{(general Neumann)}$$
for all test functions $v\in H^1(\Omega)$.
\end{enumerate}
\end{block}
\end{frame}


\begin{frame}{recall: general Dirichlet boundary conditions}

\begin{itemize}
\item general Dirichlet conditions have an additional complication
\end{itemize}

\begin{block}{weak form: general Dirichlet problem}
Suppose $f\in L^2(\Omega)$ and $g_D\in L^2(\partial\Omega)$ are given.  Let $u_D \in H^1(\Omega) \cap C(\bar\Omega)$ be any function, defined on the whole domain $\Omega$, so that $u_D=g_D$ along $\partial\Omega$.  Find a solution $u\in H_0^1(\Omega)$ so that
    $$\int_\Omega \grad u \cdot \grad v\,dx = \int_\Omega f v\,dx - \int_\Omega \grad u_D \cdot \grad v\,dx$$
for all test functions $v\in H_0^1(\Omega)$.
\end{block}

\begin{itemize}
\item the actually solution to the problem is $u+u_D$
\item even more generally, ``mixed'' conditions could partition the boundary: $\partial \Omega = \partial_D \Omega \cup \partial_N \Omega$
\end{itemize}

\end{frame}


% at start of sections 2,3,... put outline
\AtBeginSection[] {
    \begin{frame}{Outline}
    \setbeamertemplate{section in toc}[sections numbered]
    \tableofcontents[currentsection]%,hideallsubsections]
    \end{frame}
}


\section{Poisson equation on a disk (via Fourier series)}

\begin{frame}{Fourier series along the boundary of the disk}

\begin{itemize}
\item let $D = \{x\in \RR^2\,:\,|x|<1\}$ be the open unit disk, with boundary
	$$S = \partial D = \{x\in \RR^2\,:\,|x|=1\}.$$
\item suppose $g\in L^2(S,ds)$, equivalently $g\in L^2(-\pi,\pi)$
\item then from Chapter 4 we know $g$ has a complex Fourier series:
	$$g(x) = \frac{1}{\sqrt{2\pi}} \sum_{n\in\ZZ} c_n e^{inx}, \qquad c_n = \frac{1}{\sqrt{2\pi}} \int_{-\pi}^\pi g(x) e^{-inx}\,dx$$

    \begin{itemize}
    \item[$\circ$] regarding the normalization, $\{\frac{1}{\sqrt{2\pi}} e^{inx}\}_{n\in\ZZ}$ is a complete ON sequence
    \end{itemize}
\item pictures of disk and $g$:

\vspace{30mm}
\end{itemize}
\end{frame}


\begin{frame}{real Fourier series}

\begin{lemma}  The Fourier series $g(x) = \frac{1}{\sqrt{2\pi}} \sum_{n\in\ZZ} c_n e^{inx}$ is a real-valued function if and only if $c_0\in\RR$ and $c_{-n} = \overline{c_n}$.
\end{lemma}

\begin{corollary}  If $g(x)$ is a real-valued function, and if for $n\ge 0$ we write $c_n = \frac{\sqrt{2\pi}}{2} (a_n-i b_n)$, with $a_n,b_n\in\RR$, then
    $$g(x) = \frac{a_0}{2} + \sum_{n=1}^\infty a_n \cos(n\theta) + b_n \sin(n\theta).$$
\end{corollary}

\footnotesize
\noindent \emph{Proof.} Write $g(x) = \frac{1}{\sqrt{2\pi}} c_0 + \frac{1}{\sqrt{2\pi}} \sum_{n=1}^\infty \left(c_n e^{inx} + c_{-n} e^{-inx}\right)$.  Then use the lemma, giving
	$$g(x) = \frac{1}{\sqrt{2\pi}} c_0 + \frac{1}{\sqrt{2\pi}} \sum_{n=1}^\infty 2 \Re \left(c_n e^{inx}\right)$$
Substitute $c_n = \frac{\sqrt{2\pi}}{2} (a_n-i b_n)$ for $n\ge 0$, and simplify, to get the result.\qed
\normalsize
\end{frame}


\begin{frame}{harmonic functions on the disk}

\begin{itemize}
\item on the other hand, recall that the real and imaginary parts of any complex-analytic function on the disk are harmonic
\item for example, the monomials $z^n$, $n\ge 0$, are analytic
\item also, $z=r e^{i\theta}$ in polar coordinates, so the real and imaginary parts of
    $$z^n=r^n e^{in\theta}=r^n(\cos(n\theta) + i \sin(n\theta))$$
are harmonic
\end{itemize}

\begin{definition} for $n\ge 0$, define these harmonic functions on $D$:
$$\phi_n(r,\theta) = r^n \cos(n\theta), \quad \psi_n(r,\theta) = r^n \sin(n\theta)$$
\end{definition}

FIXME in polar, $\grad^2 u = u_{rr} + \frac{1}{r} u_r + \frac{1}{r^2} u_{\theta\theta}$
\end{frame}


\begin{frame}{Dirichlet's problem on a disc}

FIXME Fourier series along boundary converts to smooth and harmonic in the interior

FIXME see codes/disklaplace.m

\begin{itemize}
\item x
\end{itemize}
\end{frame}


\begin{frame}{Poisson on a disc}

FIXME perhaps google ``chebfun eigenfunctions of Laplacian on the disk''

\begin{itemize}
\item x
\end{itemize}
\end{frame}


\section{does a solution to Poisson's equation exist?}

\begin{frame}{bar}

FIXME: minimization route

\begin{itemize}
\item x
\end{itemize}
\end{frame}


\begin{frame}{bar}

FIXME: Riesz representation theorem route

\begin{itemize}
\item x
\end{itemize}
\end{frame}


\begin{frame}{bar}

FIXME: sketch how ellipticity means $u\in H^{2+k}(\Omega)$ if $f\in H^k(\Omega)$

\begin{itemize}
\item x
\end{itemize}
\end{frame}


\section{lecture content in week 13}

\begin{frame}{lecture content in week 13}

\begin{itemize}
\item most of this material is only from these slides

\begin{block}{to know from these slides:}
\begin{itemize}
\item x
\end{itemize}
\end{block}

\begin{block}{to prove:}
\begin{itemize}
\item x
\end{itemize}
\end{block}

%\item Assignment 7 is posted at \quad  \href{https://bueler.github.io/fa/index.html}{\texttt{bueler.github.io/fa}}
\end{itemize}
\end{frame}

\end{document}
